\begin{proposition}[twist independence in the sphere cases]\label{push:src:C.4.1}
For the fixed multiplicities $(3,4)$ and all admissible data in Paper~I, Proposition~4.5 and Paper~I, Proposition~4.6 and Paper~I, Proposition~4.7 and Paper~I, Proposition~4.8
with $|p|=1$, the monodromy class of the positive $w$-circle is the
same class in $\Theta_7$, with the complex orientation on each fibre.
\end{proposition}
\Uses{the hypotheses, notation and definitions stated here; \cref{push:src:C.3.24}}


\begin{proof}
The explicit sufficiency maps in E.1 depend holomorphically on a
small $c$-disc: the lattice defects and toric integral exponents are
constant, the periods and unit factors are holomorphic, and the
translation correction is solved in the fixed sheaf $V$ with
$H^1(B,V)=0$. This is the parameter-dependent assertion of Paper~I, Corollary~23.3,
not just the fibrewise classification. For the second twist
coefficient $s_2$, the parameter identification is
$W=e^{-\pi i(c+s_2)}=(-1)^{s_2}w$.
The rotation of the punctured disc preserves its positive
fundamental-group generator. Local family identifications differ
by elements of $\Aut(X)=\C^*$, a connected group, so their
transition changes are smoothly isotopic to the identity. Following
these identifications around the circle conjugates the monodromy
in the oriented mapping-class group. Since $\Theta_7$ is abelian,
the resulting class is unchanged. If the lattice sign in Paper~I, Proposition~4.5 and Paper~I, Proposition~4.6 and Paper~I, Proposition~4.7 and Paper~I, Proposition~4.8 is
$-1$, the comparison map is nevertheless holomorphic and preserves
the complex orientation; no extra sign is introduced.
The common class is zero by \ref{push:src:C.3.24}. This last numerical consequence
uses only the already proved equality of the classes, not an
additional deformation or conjugation argument.
\end{proof}

